Distinguishing central dark matter cores from cusps in dwarf galaxies requires inference that accounts for limited kinematic information and prior dependence. We test the sensitivity of core--cusp inference to a joint Jeffreys prior, defined by the expected Fisher information of a conditional Gaussian stellar-velocity likelihood, relative to a prior uniform in specified sampling coordinates. JeansPy supplies Jeans predictions and automatic parameter derivatives, including systemic-velocity information. With common seven-parameter support, we examine cored and NFW distribution-function catalogues in two settings: 256-star Draco-like and 64-star Segue 1-like mocks. Systemic velocity is analytically integrated while retaining its joint information factor. Independent direct ensemble sampling passes the frozen sampling requirements for all eight targets; selected-point numerical checks and a follow-up precision audit support the plotted comparisons. For Draco-like catalogues, changing the prior shifts the median inner slope from 0.856 to 0.666 for the NFW truth of 1, and from 0.879 to 0.460 for the core truth of 0. The Segue-like medians change from 0.688 to 0.509 and from 0.412 to 0.409, respectively. Broad intervals overlap, and the nearly unchanged Segue core slope coexists with substantial changes in inferred density. These fixed realizations demonstrate prior sensitivity, not improved discrimination or resolved likelihood degeneracies. Repeated likelihood-matched and dynamical mocks remain necessary to assess recovery and coverage.
